\documentclass[11pt,a4paper]{article}

% --- Core Typography & Geometry ---
\usepackage[utf8]{inputenc}
\usepackage[T1]{fontenc}
\usepackage{lmodern}
\usepackage[margin=0.85in, top=0.85in, bottom=0.85in]{geometry}
\usepackage{amsmath,amssymb,amsthm,mathtools}
\usepackage{enumitem}
\usepackage{graphicx}
\usepackage{xcolor}
\usepackage{microtype}
\usepackage[colorlinks=true,linkcolor=utblue,citecolor=utblue,urlcolor=utblue]{hyperref}

% --- University Color Palette ---
\definecolor{utblue}{RGB}{12, 62, 116}        % Official Tehran University Blue (#0C3E74)
\definecolor{darktext}{RGB}{15, 23, 42}       % Dark Slate Text

% --- Paragraph and Spacing Settings ---
\setlength{\parindent}{0pt}
\setlength{\parskip}{5pt}

% --- Problem Header Macro ---
\newcommand{\problemheader}[2]{%
  \noindent
  \begin{minipage}{\textwidth}
    {\large\bfseries\textcolor{utblue}{#1}}\hfill
    {\large\bfseries\textcolor{utblue}{#2}}\\[3pt]
    \textcolor{utblue}{\rule{\textwidth}{1pt}}
  \end{minipage}%
  \par\vspace{8pt}
}

\begin{document}

% --- Header with Dual Logos ---
\noindent
\begin{minipage}[c]{0.13\textwidth}
  \includegraphics[width=\linewidth, keepaspectratio]{../../assets/images/ut-logo.png}
\end{minipage}%
\hfill
\begin{minipage}[c]{0.70\textwidth}
  \centering
  {\Large\bfseries\sffamily\textcolor{utblue}{Basics of Data Science}}\\[2pt]
  {\small\textbf{Fall 2026}}\\[3pt]
  {\small\textbf{Instructor:} Sayeh Mirzaee}\\[1.5pt]
  {\small\textbf{TAs:} Iliya Jafarizadeh}\\[1.5pt]
  {\small\textbf{Deadline:} October 8, 2026 (23:59)}\\[3pt]
  {\large\bfseries\textcolor{darktext}{Homework 0}}\\[3pt]
  {\footnotesize\textbf{Topics:} Supervised Learning Foundations, Bias-Variance Decomposition, Augmented Ridge Regression, Softmax Concavity, and Cross-Validation Leakage}
\end{minipage}%
\hfill
\begin{minipage}[c]{0.15\textwidth}
  \raggedleft
  \includegraphics[width=\linewidth, keepaspectratio]{../../assets/images/faculty-logo.png}
\end{minipage}

\vspace{6pt}
\noindent\textcolor{utblue}{\rule{\textwidth}{1.3pt}}
\vspace{12pt}

% ==============================================================================
% PROBLEM 1
% ==============================================================================
\problemheader{Problem 1: Expected Prediction Error \& Bias-Variance in $k$-NN}{[25 points]}

Consider the classical nonparametric regression setting. Suppose we have a training dataset $\mathcal{D} = \{(x_i, y_i)\}_{i=1}^N$ where the scalar targets are generated by the model:
\[
  y = f(x) + \epsilon, \quad \text{with } \mathbb{E}[\epsilon \mid x] = 0 \text{ and } \operatorname{Var}(\epsilon \mid x) = \sigma^2,
\]
where the additive noise $\epsilon$ is independent across observations and independent of the input vector $x \in \mathbb{R}^p$.

For a fixed query test point $x_0 \in \mathbb{R}^p$, consider estimating the unknown regression function $f(x_0)$ using the $1$-Nearest-Neighbor ($1$-NN) estimator:
\[
  \hat{f}_1(x_0) = y_{(1)},
\]
where $x_{(1)} \in \{x_1, \dots, x_N\}$ denotes the training input closest to $x_0$ under Euclidean distance, and $y_{(1)}$ is its associated observed target.

\begin{enumerate}[label=(\alph*), itemsep=5pt, leftmargin=18pt]
  \item \textbf{[9 points]} Under the standard squared-error loss function $L(y, \hat{f}) = (y - \hat{f})^2$, show that the conditional Expected Prediction Error at $x_0$, evaluated with respect to a new independent target observation $y_0 = f(x_0) + \epsilon_0$ and conditioned on the training inputs $\{x_i\}_{i=1}^N$, decomposes exactly as:
  \[
    \operatorname{EPE}(x_0) \equiv \mathbb{E}_{\mathcal{D}, \epsilon_0} \left[ (y_0 - \hat{f}_1(x_0))^2 \;\middle|\; x_0, \{x_i\}_{i=1}^N \right] = \sigma^2 + \left[ f(x_0) - f(x_{(1)}) \right]^2 + \sigma^2.
  \]
  Explicitly identify each of the three terms in this decomposition (the irreducible error, the squared model bias, and the estimator sampling variance).

  \item \textbf{[9 points]} Generalize this result to the $k$-Nearest-Neighbor ($k$-NN) averaging estimator:
  \[
    \hat{f}_k(x_0) = \frac{1}{k} \sum_{\ell=1}^k y_{(\ell)},
  \]
  where $x_{(1)}, \dots, x_{(k)}$ are the $k$ closest training points to $x_0$. Derive the explicit formula for $\operatorname{EPE}(x_0; k)$ as a function of $k$, the underlying function values $\{f(x_{(\ell)})\}_{\ell=1}^k$, and the noise variance $\sigma^2$.

  \item \textbf{[7 points]} Analyze the behavior of $\operatorname{EPE}(x_0; k)$ as $k$ varies from $1$ to $N$:
  \begin{itemize}[itemsep=2pt, leftmargin=14pt]
    \item Explain mathematically what happens to the bias term and the variance term when $k = 1$ versus $k = N$.
    \item Relate this behavior to the fundamental concepts of model complexity, underfitting, and overfitting.
    \item Briefly discuss why the bias term fails to vanish in high dimensions ($p \gg 1$) even for $k=1$ (the curse of dimensionality).
  \end{itemize}
\end{enumerate}

\newpage

% ==============================================================================
% PROBLEM 2
% ==============================================================================
\problemheader{Problem 2: Ridge Regression as an Augmented OLS System}{[25 points]}

In high-dimensional and collinear data regimes, Ordinary Least Squares can suffer from severe variance inflation or non-invertibility of $X^T X$. Ridge regression stabilizes parameter estimates by penalizing the $L_2$-norm of the regression coefficients:
\[
  \min_{\beta \in \mathbb{R}^p} \left\{ \|y - X\beta\|_2^2 + \lambda \|\beta\|_2^2 \right\},
\]
where $\lambda > 0$ is a fixed regularization hyperparameter.

\begin{enumerate}[label=(\alph*), itemsep=5pt, leftmargin=18pt]
  \item \textbf{[5 points]} Define an augmented feature matrix $X^* \in \mathbb{R}^{(N+p) \times p}$ and an augmented response vector $y^* \in \mathbb{R}^{N+p}$ by vertically stacking:
  \[
    X^* = \begin{bmatrix} X \\ \sqrt{\lambda} I_p \end{bmatrix}, \quad y^* = \begin{bmatrix} y \\ 0_p \end{bmatrix},
  \]
  where $I_p$ is the $p \times p$ identity matrix and $0_p$ is a $p$-dimensional vector of zeros.

  \item \textbf{[6 points]} Consider setting up an ordinary (unregularized) least squares problem on this augmented data:
  \[
    \min_{\beta \in \mathbb{R}^p} \|y^* - X^* \beta\|_2^2.
  \]
  Show algebraically that this augmented OLS objective is identically equal to the Ridge regression objective function.

  \item \textbf{[7 points]} Show that the closed-form OLS solution on the augmented dataset:
  \[
    \hat{\beta}^* = (X^{*T} X^*)^{-1} X^{*T} y^*
  \]
  reproduces the regularized Ridge regression estimator:
  \[
    \hat{\beta}_{\text{ridge}} = (X^T X + \lambda I_p)^{-1} X^T y.
  \]

  \item \textbf{[7 points]} Prove that for any $\lambda > 0$, the matrix $X^{*T} X^*$ is strictly positive definite:
  \[
    v^T (X^{*T} X^*) v > 0 \quad \forall v \in \mathbb{R}^p \setminus \{0_p\},
  \]
  regardless of whether $N < p$ or whether the original matrix $X$ has linearly dependent columns. Explain how this algebraic property guarantees that the Ridge normal equations always possess a unique and numerically stable solution.
\end{enumerate}

\newpage

% ==============================================================================
% PROBLEM 3
% ==============================================================================
\problemheader{Problem 3: Multiclass Softmax Regression \& Log-Likelihood Concavity}{[25 points]}

Consider a classification problem with $K \ge 2$ mutually exclusive classes. For each sample $i \in \{1, \dots, N\}$, the response $y_i \in \{1, \dots, K\}$ is encoded via a 1-of-$K$ indicator vector $y_i = [y_{i1}, \dots, y_{iK}]^T \in \{0, 1\}^K$ such that $\sum_{k=1}^K y_{ik} = 1$.

The multiclass logistic regression model (Softmax classifier) models class posterior probabilities as:
\[
  p_k(x_i; \theta) \equiv P(Y = k \mid x_i) = \frac{\exp(\beta_k^T x_i)}{\sum_{\ell=1}^K \exp(\beta_\ell^T x_i)}, \quad k = 1, \dots, K,
\]
where $x_i \in \mathbb{R}^{p+1}$ (including the constant bias feature $1$) and $\theta = [\beta_1^T, \dots, \beta_K^T]^T \in \mathbb{R}^{K(p+1)}$ gathers all parameter vectors.

\begin{enumerate}[label=(\alph*), itemsep=5pt, leftmargin=18pt]
  \item \textbf{[5 points]} Formulate the conditional log-likelihood function $\ell(\theta)$ for an independently sampled training dataset $\mathcal{D} = \{(x_i, y_i)\}_{i=1}^N$:
  \[
    \ell(\theta) = \sum_{i=1}^N \sum_{k=1}^K y_{ik} \log p_k(x_i; \theta).
  \]

  \item \textbf{[7 points]} Derive the gradient vector of the log-likelihood with respect to an arbitrary class weight vector $\beta_k$ ($k \in \{1, \dots, K\}$), showing that:
  \[
    \nabla_{\beta_k} \ell(\theta) \equiv \frac{\partial \ell(\theta)}{\partial \beta_k} = \sum_{i=1}^N \left( y_{ik} - p_k(x_i; \theta) \right) x_i.
  \]
  Provide an intuitive physical interpretation of this gradient update rule in terms of prediction error and feature weighting.

  \item \textbf{[6 points]} Compute the block components of the Hessian matrix $\mathcal{H}_{k\ell} \equiv \frac{\partial^2 \ell(\theta)}{\partial \beta_k \partial \beta_\ell^T} \in \mathbb{R}^{(p+1) \times (p+1)}$, proving that:
  \[
    \frac{\partial^2 \ell(\theta)}{\partial \beta_k \partial \beta_\ell^T} = -\sum_{i=1}^N p_k(x_i; \theta) \left( \delta_{k\ell} - p_\ell(x_i; \theta) \right) x_i x_i^T,
  \]
  where $\delta_{k\ell}$ is the Kronecker delta ($\delta_{k\ell} = 1$ if $k = \ell$, and $0$ otherwise).

  \item \textbf{[7 points]} Prove that the negative Hessian $-\mathcal{H} = -\nabla^2 \ell(\theta)$ is positive semi-definite everywhere in parameter space:
  \[
    u^T (-\mathcal{H}) u \ge 0 \quad \forall u \in \mathbb{R}^{K(p+1)}.
  \]
  Conclude that $\ell(\theta)$ is strictly concave (up to an invariant additive constant), guaranteeing that gradient ascent and Newton-Raphson methods converge to a unique global optimum.
\end{enumerate}

\newpage

% ==============================================================================
% PROBLEM 4
% ==============================================================================
\problemheader{Problem 4: Feature Selection \& Data Leakage in Cross-Validation}{[25 points]}

Data leakage during model validation is one of the most pervasive failure modes in machine learning and empirical data science.

Consider the following simulation experiment:
\begin{itemize}[itemsep=2pt, leftmargin=14pt]
  \item A synthetic dataset contains $N = 50$ samples and $p = 5000$ predictors.
  \item All feature entries $X_{ij} \overset{\text{i.i.d.}}{\sim} \mathcal{N}(0, 1)$ ($i = 1, \dots, 50$; $j = 1, \dots, 5000$) are sampled completely independently of each other.
  \item The binary class labels $Y_i \in \{0, 1\}$ are assigned randomly by tossing an unbiased coin:
  \[
    P(Y_i = 1) = P(Y_i = 0) = 0.5.
  \]
\end{itemize}
Because the response $Y$ and all predictors $X$ are completely independent, the true population Bayes optimal classifier achieves a generalization error rate of exactly $\operatorname{Err}_{\text{true}} = 50\%$.

\begin{enumerate}[label=(\alph*), itemsep=5pt, leftmargin=18pt]
  \item \textbf{[9 points] (The Flawed Protocol: Data Leakage)} \\
  A practitioner executes the following validation procedure:
  \begin{enumerate}[label=\arabic*., itemsep=2pt]
    \item Across the entire dataset of $N = 50$ samples, compute the sample correlation $r_j = \widehat{\operatorname{Corr}}(X_j, Y)$ for each of the $p = 5000$ features.
    \item Select the subset $\mathcal{S} \subset \{1, \dots, p\}$ containing the $m = 100$ predictors with the largest absolute sample correlation $|r_j|$.
    \item Discard the other $4900$ features and evaluate a $1$-Nearest-Neighbor or Linear Discriminant classifier using standard $10$-fold cross-validation restricted to the $100$ features in $\mathcal{S}$.
  \end{enumerate}
  Explain mathematically and conceptually why this procedure will report a $10$-fold cross-validation error rate near $0\%$, dramatically misleading the researcher into believing they have discovered a near-perfect predictive model.

  \item \textbf{[7 points] (Extreme Value Statistics of Spurious Correlation)} \\
  When $N = 50$, what is the distribution of the sample correlation $r_j$ under the null hypothesis of independence? \\
  Given that $p = 5000 \gg N$, analyze the expected behavior of the extreme order statistic:
  \[
    \mathbb{E}\left[ \max_{1 \le j \le 5000} |r_j| \right].
  \]
  Show why having a massive search space of features guarantees that several dozen features will exhibit strong empirical correlation with $Y$ purely by random chance.

  \item \textbf{[9 points] (The Correct Protocol)} \\
  Detail the precise, algorithmic protocol for conducting $k$-fold cross-validation correctly when supervised feature screening is required. Explicitly describe how the $100$ features must be selected within each fold, and explain why this eliminates all information leakage between training and testing splits.
\end{enumerate}

\end{document}
